\documentclass[12pt]{article}

% ----
% Packages
% ----
\usepackage{fontspec}
\usepackage{amsmath}
\setmainfont{Arial}
\usepackage{pgfplots}
\pgfplotsset{compat=1.18}
\usepgfplotslibrary{dateplot}
\usepgfplotslibrary{groupplots}
\usepackage{pgfplotstable}
\usepackage{booktabs}
\usepackage{placeins}   % \FloatBarrier: keep tables and figures in source order

\usepackage[left=3cm, right=3cm, top=3cm, bottom=3cm]{geometry}





% ----
% Read Files
% ----
% Crisis excess return vs. MCap monthly, 19 portfolios (final_data_prep.ipynb)
\newcommand{\scatterfile}{data_finalized/figure_scatter_frequency.csv}
% Monthly p_crisis of the three detectors (final_data_prep.ipynb)
\newcommand{\pcrisismonthlyfile}{data_finalized/figure_pcrisis_monthly.csv}
% Daily value of the seven monthly HRP and LSTM portfolios (final_data_prep.ipynb)
\newcommand{\valuefile}{data_finalized/figure_value_hrp_lstm.csv}
% Cumulative log return of MCap and the seven LSTM portfolios per crisis window,
% one file per crisis: <prefix>1.csv ... <prefix>7.csv (final_data_prep.ipynb)
\newcommand{\crisespanelfile}{data_finalized/figure_crises_lstm_c}
% Tables (final_data_prep.ipynb)
\newcommand{\tableoverviewfile}{data_finalized/table_1_crisis_overview.csv}
\newcommand{\tablefullperiodfile}{data_finalized/table_2_full_period.csv}
\newcommand{\tablecrisisavgfile}{data_finalized/table_3_crisis_average.csv}
\newcommand{\tabledetectorsfile}{data_finalized/table_4_detectors.csv}
\newcommand{\tabledetectorinfluencefile}{data_finalized/table_6_detector_influence.csv}

% Model family colours (Okabe-Ito, colour-blind safe); use the same in all charts
\definecolor{famBenchmark}{RGB}{110,110,110}   % market_cap
\definecolor{famSimple}{RGB}{230,159,0}        % equal_weight, mvo
\definecolor{famClassical}{RGB}{0,158,115}          % hrp
\definecolor{famML}{RGB}{0,114,178}            % xgb, rf, lstm
% One colour per model (Okabe-Ito, colour-blind safe), for charts that show
% every model separately
\definecolor{modMCap}{RGB}{110,110,110}
\definecolor{modEqual}{RGB}{230,159,0}       % 1/N: orange
\definecolor{modHRP}{RGB}{0,158,115}         % bluish green
\definecolor{modMVO}{RGB}{213,94,0}          % vermillion
\definecolor{modLSTM}{RGB}{0,114,178}        % blue
\definecolor{modRF}{RGB}{86,180,233}         % sky blue
\definecolor{modXGB}{RGB}{204,121,167}       % reddish purple
% Label anchor opposite to the side the label should sit on (label_pos = west -> anchor east)
\pgfkeys{/labelanchor/.cd, west/.initial=east, east/.initial=west,
         north/.initial=south, south/.initial=north}
\newcommand{\labelanchor}[1]{\pgfkeysvalueof{/labelanchor/#1}}

% Time-series line chart 1998-2025, shared by all figures
\pgfplotsset{timechart/.style={
    width=\textwidth, height=5.5cm,
    date coordinates in=x,
    xmin=1998-01-01, xmax=2025-12-31,
    xtick={2000-01-01,2005-01-01,2010-01-01,2015-01-01,2020-01-01,2025-01-01},
    restrict x to domain*=1998-01-01:2025-12-31,
    xticklabel=\year,
    grid=major, grid style={gray!20},
    tick label style={font=\small},
    table/col sep=comma,
}}

% p_crisis line chart, shared by all three regime figures
\pgfplotsset{pcrisis/.style={
    timechart,
    ymin=0, ymax=1,
    ytick={0,0.5,1},
    ylabel={$p_{\mathrm{crisis}}$},
}}

% S&P 500 cumulative log return underlay (right axis), drawn below a pcrisis axis
\newcommand{\spfile}{sp500_daily_returns.csv}
\pgfplotsset{spunderlay/.style={
    timechart,
    ymin=-0.5, ymax=2,
    ytick={-0.5,0,0.5,1,1.5,2},
    axis y line*=right, axis x line=none,
    ylabel={S\&P 500 cum.\ log return},
    grid=none,
}}
\newcommand{\regimepanelheight}{3cm}
% One panel of the monthly stacked figure: #1 = column of \pcrisismonthlyfile,
% #2 = y label, #3 = axis options.  Step line: the value of the first trading
% day of a month is held until the next month.
\newcommand{\pcrisispanel}[3]{%
    \renewcommand{\crisisymin}{-0.06}%
    \renewcommand{\crisisymax}{1.06}%
    \begin{axis}[timechart, scale only axis, width=0.88\textwidth, height=\regimepanelheight,
                 ymin=-0.06, ymax=1.06, ytick={0,0.5,1}, ylabel={#2}, axis on top,
                 ymajorgrids=false, #3]
        \crises
        \addplot[forget plot, black!55, dashed, line width=0.4pt] coordinates {(1998-01-01,0.5) (2025-12-31,0.5)};
        \addplot[const plot, blue!70!black, line width=0.6pt, mark=none] table[x=date, y=#1] {\pcrisismonthlyfile};
    \end{axis}%
}

% Regime portfolio colours: one Okabe-Ito hue per portfolio (colour-blind safe),
% the three cash and the three alternative hues far apart from each other;
% cash also dashed, alternative solid.  The same in every chart.
\definecolor{regBase}{RGB}{0,0,0}
\definecolor{regCPcash}{RGB}{230,159,0}        % changepoint cash: orange
\definecolor{regCLcash}{RGB}{86,180,233}       % clustering cash: sky blue
\definecolor{regHMMcash}{RGB}{204,121,167}     % HMM cash: reddish purple
\definecolor{regCP}{RGB}{213,94,0}             % changepoint alternative: vermillion
\definecolor{regCL}{RGB}{0,114,178}            % clustering alternative: blue
\definecolor{regHMM}{RGB}{0,158,115}           % HMM alternative: bluish green
\pgfplotsset{
    regbase/.style={regBase, line width=0.6pt},
    regcash/.style={#1, line width=0.6pt, dash pattern=on 2.5pt off 1pt},
    regalt/.style={#1, line width=0.6pt},
}

% One panel of the stacked value figure: the seven monthly portfolios of one
% model.  #1 = column prefix in \valuefile (hrp, lstm), #2 = y label,
% #3 = axis options (position, legend, x tick labels), #4 = extra content
% (e.g. \crisislabels).  Log scale, value of 1 invested at the start.  Crises:
% shading and the trough line only.
\newcommand{\valuepanelheight}{4.5cm}
\newcommand{\valuepanel}[4]{%
    \renewcommand{\crisisymin}{0.7}%
    \renewcommand{\crisisymax}{25}%
    \begin{axis}[timechart, scale only axis, width=0.88\textwidth, height=\valuepanelheight,
                 ymode=log, ymin=0.7, ymax=25, ytick={1,2,5,10,20}, yticklabels={1,2,5,10,20},
                 ylabel={#2}, axis on top, ymajorgrids=false, clip mode=individual, #3]
        \crisesshade
        \crisislist\crisistrough
        #4
        % legend order (4 columns): Baseline, cash x 3 / empty, alternative x 3;
        % the baseline is drawn last so it stays visible on top
        \addlegendimage{regbase}
        \addplot[regcash=regCPcash]  table[x=date, y=#1_changepoint_cash] {\valuefile};
        \addplot[regcash=regCLcash]  table[x=date, y=#1_clustering_cash]  {\valuefile};
        \addplot[regcash=regHMMcash] table[x=date, y=#1_hmm_cash]         {\valuefile};
        \addlegendimage{empty legend}
        \addplot[regalt=regCP]       table[x=date, y=#1_changepoint_alt]  {\valuefile};
        \addplot[regalt=regCL]       table[x=date, y=#1_clustering_alt]   {\valuefile};
        \addplot[regalt=regHMM]      table[x=date, y=#1_hmm_alt]          {\valuefile};
        \addplot[regbase, forget plot] table[x=date, y=#1_baseline]       {\valuefile};
    \end{axis}%
}

% Crisis panels (4 x 2 groupplot): full calendar months around one crisis,
% cumulative log return relative to the S&P peak (weekly points); one file per
% crisis (\crisespanelfile<k>.csv), so nothing has to be filtered.
\newcommand{\shortmonth}[1]{\ifcase#1\or Jan\or Feb\or Mar\or Apr\or May\or Jun\or Jul\or Aug\or Sep\or Oct\or Nov\or Dec\fi}
% #1 = crisis number, #2 = title, #3 = 1st of the peak month, #4 = peak,
% #5 = trough, #6 = even, #7 = last day of the even month (dates).
% Inside a groupplot.  Light grey shading peak -> even; lines: MCap (thin
% grey), the six regime variants (colours as \valuepanel), the LSTM baseline
% (thick black, drawn last); dashed vertical line at the S&P trough.
\newcommand{\crisispanel}[7]{%
    \nextgroupplot[title={#2}, date coordinates in=x, xmin=#3, xmax=#7, xtick={#4,#5,#6},
                   xticklabel={\shortmonth{\month}~\year}, table/col sep=comma,
                   every axis plot/.append style={mark=none}]
        \fill[gray!15] ({axis cs:#4,0} |- {rel axis cs:0,0}) rectangle ({axis cs:#6,0} |- {rel axis cs:0,1});
        \addplot[gray!60, line width=0.5pt]  table[x=date, y=mcap]                  {\crisespanelfile#1.csv};
        \addplot[regcash=regCPcash]          table[x=date, y=lstm_changepoint_cash] {\crisespanelfile#1.csv};
        \addplot[regcash=regCLcash]          table[x=date, y=lstm_clustering_cash]  {\crisespanelfile#1.csv};
        \addplot[regcash=regHMMcash]         table[x=date, y=lstm_hmm_cash]         {\crisespanelfile#1.csv};
        \addplot[regalt=regCP]               table[x=date, y=lstm_changepoint_alt]  {\crisespanelfile#1.csv};
        \addplot[regalt=regCL]               table[x=date, y=lstm_clustering_alt]   {\crisespanelfile#1.csv};
        \addplot[regalt=regHMM]              table[x=date, y=lstm_hmm_alt]          {\crisespanelfile#1.csv};
        \addplot[black, line width=1.1pt]    table[x=date, y=lstm_baseline]         {\crisespanelfile#1.csv};
        \draw[black!60, dashed, line width=0.5pt]
            ({axis cs:#5,0} |- {rel axis cs:0,0}) -- ({axis cs:#5,0} |- {rel axis cs:0,1});
}

% Portfolio value (NAV) line chart
\pgfplotsset{navchart/.style={
    timechart,
    ymin=0.5, ymax=4,
    ytick={1,2,3,4},
    ylabel={NAV},
    axis on top,
}}

% Crisis periods (crises.txt): grey shading from peak to even,
% vertical red lines at peak, trough and even.
% Vertical extent is \crisisymin to \crisisymax (set to the axis ymin/ymax).
% \crisesshade and \criseslines draw the two parts separately, so other
% layers can go between the shading and the lines.
% Colours of the shading and the lines; renew them inside a figure to change one chart only.
\newcommand{\crisisymin}{0}
\newcommand{\crisisymax}{1}
\newcommand{\crisisshadecolor}{gray!25}
\newcommand{\crisislinecolor}{red}
\newcommand{\crisisshade}[3]{%
    \addplot[forget plot, draw=none, fill=\crisisshadecolor] coordinates {(#1,\crisisymin) (#3,\crisisymin) (#3,\crisisymax) (#1,\crisisymax)} --cycle;
}
\newcommand{\crisismarks}[3]{%
    \addplot[forget plot, \crisislinecolor, line width=0.4pt, mark=none] coordinates {(#1,\crisisymin) (#1,\crisisymax)};
    \addplot[forget plot, \crisislinecolor, line width=0.4pt, mark=none] coordinates {(#2,\crisisymin) (#2,\crisisymax)};
    \addplot[forget plot, \crisislinecolor, line width=0.4pt, mark=none] coordinates {(#3,\crisisymin) (#3,\crisisymax)};
}
\newcommand{\crisis}[3]{\crisisshade{#1}{#2}{#3}\crisismarks{#1}{#2}{#3}}
% #1 = macro applied to every crisis (peak, trough, even)
\newcommand{\crisislist}[1]{%
    #1{1998-07-17}{1998-08-31}{1998-11-23}% C1
    #1{2000-03-24}{2002-10-09}{2007-05-30}% C2
    #1{2007-10-09}{2009-03-09}{2013-03-28}% C3
    #1{2018-09-20}{2018-12-24}{2019-04-23}% C4
    #1{2020-02-19}{2020-03-23}{2020-08-18}% C5
    #1{2022-01-03}{2022-10-12}{2024-01-19}% C6
    #1{2025-02-19}{2025-04-08}{2025-06-27}% C7
}
% Trough line only, and the labels C1 ... C7 above the axis at the trough
\newcommand{\crisistrough}[3]{%
    \addplot[forget plot, \crisislinecolor, line width=0.4pt, mark=none] coordinates {(#2,\crisisymin) (#2,\crisisymax)};
}
\newcounter{crisisno}
\newcommand{\crisislabel}[3]{%
    \stepcounter{crisisno}%
    \edef\crisislabelnode{\noexpand\node[anchor=south, inner sep=1.5pt, font=\noexpand\scriptsize]
        at (axis cs:#2,\crisisymax) {C\arabic{crisisno}};}%
    \crisislabelnode
}
\newcommand{\crisislabels}{\setcounter{crisisno}{0}\crisislist\crisislabel}
\newcommand{\crises}{\crisislist\crisis}
\newcommand{\crisesshade}{\crisislist\crisisshade}
\newcommand{\criseslines}{\crisislist\crisismarks}

% ----
% Calculations
% ----





% ----
% Begin Doc
% ----

\begin{document}

\begin{titlepage}
    \centering
    \textbf{Thesis -> Paper}\vspace{3cm}\\
    \textbf{Benedikt Löw}\\
\end{titlepage}






% ----
% Abstract
% ----

\newpage
\pagestyle{plain}
\pagenumbering{Roman}
\setcounter{page}{2}
\begin{center}
    \textbf{Abstract}
\end{center}

\noindent

This thesis is being transformed into a paper


% ----
% Main Part
% ----

\newpage
\pagestyle{plain}
\pagenumbering{arabic}
\setcounter{page}{1}
\section{Introduction}
\subsection{Motivation and Research Context}
The mean-variance framework introduced by Markowitz [1952] has shaped 
capital allocation for over seven decades, yet its core assumptions of 
stable correlations and reliable parameter estimates break down in exactly 
the conditions where diversification is meant to help. During market declines, 
correlations among S\&P 500 constituents rise well above their upside levels, 
an asymmetry that erodes the diversification the framework is built to deliver 
precisely when protection matters most [Ang and Chen 2002].

\section{Research Question}
To what extent do classical and machine learning portfolio construction approaches deliver risk-adjusted 
outperformance and drawdown resilience across crisis episodes, and can regime detection improve them?\\

Sub questions:\\
1. How do MVO, HRP, XGB, RF and LSTM perform against market-cap weighted benchmarks, over the full sample 
and within crisis episodes?\\
2. Are the results robust to monthly, quarterly and yearly rebalancing?\\
3. How accurately and how early do changepoint, HMM and clustering detectors identify crisis episodes?\\
4. Does moving to T-bills in detected crises improve each approach, and is any gain due to timing rather 
than lower average equity exposure?\\
5. Does a regime-switched combination (ML in calm, HRP in crisis) outperform each approach alone and a 
static mix of both?






\section{Methodology}
\subsection{Time Frame}
- 28 years, January 1998 untill December 2025
\subsection{Rebalancing}
- Rebalancing Schedule 3x, monthly, quarterly, yearly
- if non trading day, next day in the month/year
\subsection{Weight Constraints}
- 1\% min and 10\% maximum allocation per stock per rebalancing date. 
- Intra rebalancing drift can exceed 10\%
\subsection{Lookback Windows}
- 60 month calculation window for HRP and MVO
- Fixed 60 month lookback window for all Machine Learning models | 24 month validation | + 1, 3, 12 month embargo
\subsection{Coverage and Missing Data}
- price data of min. 50\% given, otherwise excluded
- 5 days forward filled, otherwise missing values
\subsection{Transaction Cost}
- 0bps, 10bps, 20bps calculated. 10bps standard transaction cost
- Turnover is getting calculated and reported
\subsection{Point-in-Time}
- Everything is point in time only 
\subsection{Features}
- 

\section{Regime Implementation}
\subsection{Changepoint Detection}
\subsubsection{Step 1}
- daily VIX log changes 
- k-means clustering and mood rank test. is the spread of teh data after k different 
from before k? -> keeps the largest test statistic
- if that maximum exceeds a threshold, it raises an alarm on day t (detection date). the k with the largest
statistic becomes teh estimated break date t, search restarts from there
- ARL0 = 10000
- Mood test uses ranks, it assumes no distribution and it reacts to changes in volatility
\subsubsection{Step 2}
- estimate teh volatility once
- each detection day t, an EWMA ($\lambda$ = 0.95) of squared S\&P 500 returns from t det to t gives
the new regime's annualised volatility o. Always SPY returns, because 20\% annual volatility is from there
- Freeze it. o stays constant until the new detection. That is why your chart has steps and late exits: 
calm returns only after a new break is detected and the new segment’s vol is below 20\%.
- Map it to a score: $p_{\text{crisis}} = \Phi((\sigma - 0.20) / 0.05)$, where $\Phi$ is the standard normal CDF. So o = 15\%
gives 0.16, 20\% gives 0.50, 25\% gives 0.84 and etc.
- shift by one day. the value on day d uses data only through d-1, which mathces the paper's one day 
implementation delay

\subsection{Hidden Markov Model}
\subsubsection{Step 1}
- daily log return of S\&P 500
- states follow a markov chain. false L holds the probabilities of staying in each state
- never the state, only returns. the model infers the state from how well each day's return fits each 
state's distribution
- crisis is defined as teh state with teh hgher o. 
\subsubsection{Step 2}
- on the first 260 days, the parameters mu, o, false L are estimated by maximum likelihood
\subsubsection{Step 3}
- update parameters daily
- Each new return moves the parameters one small step toward what fits recent data best.
- Older days are down-weighted by $\lambda = 1 - 1/260$, which gives about a one-year memory.
- So $\sigma_{\text{calm}}$ and $\sigma_{\text{crisis}}$ drift over time. “Crisis” means high volatility relative to the last year or so, not above a fixed 20\%.
\subsubsection{Step 4}
- $\alpha_t \propto \alpha_{t-1} \cdot \Gamma \cdot P(r_t)$. This is Bayes’ rule:
- Prior: yesterday’s state probabilities, pushed through $\Gamma$ (regimes tend to persist).
- Times the likelihood: how probable today’s return is under each state’s normal density.
- Example: a −3\% day is about 4.3 standard deviations under calm but only 1.9 under crisis. That makes it 
roughly 900 times more likely under crisis, so the probability jumps toward crisis. One quiet day after that 
barely moves it back, because $\Gamma$ says crises persist.
\subsubsection{Step 5}
- $P(S_{t+1}) = \alpha_t \cdot \Gamma$. forecast = $p_{\text{crisis}}$
- It is made at the close of d−1 and stored on day d, the same one-day delay as the changepoint files.
- MVO and HRP switch to min-variance when it exceeds 0.5. ML uses the value directly.

\subsection{Clustering (Wasserstein K-Means)}
Wasserstein k-means compares whole monthly return distributions. It sorts past months into two groups by shape, 
then asks which group the latest month looks more like. $p_{\text{crisis}}$ is a relative distance to the two group centres, so 
it is a score, not a model probability.
\subsubsection{Step 1}
- turn each month into a point clodu
- a window of 21 daily returns becomes a cloud of 21 points. In the main run each point is (S\&P log return, VIX
log change)
- each series is divided by its expanding standard deviation up to the day before the window starts. Without this, 
VIX changes about 4x as volatile would dominate. Only past data is used
- The cloud captures the whole joint distribution: spread, fat tails, skew and how the S\&P 500 and VIX move together.
A single volatility number captures only the first of these
\subsubsection{Step 2}
- measure how different the two months are
- The wasserstein distnace W2 mathces each point of one cloud to one point of the other so that the total squared 
distance is as small as possible (hungarian algorithm). $W_2^2 = \frac{1}{21} \min_\pi \sum_i \lVert X_i - Y_{\pi(i)} \rVert^2$. 
Two calm months are close. A calm and a crisis month are far apart becasue the crisis cloud is spread wider. 
\subsubsection{Step 3}
- On the first trading day of each month, k-means with K=2 runs on all windows that end before that day. Windows slide 
by 1 day, and the history expands
- Each cluster's centre is a medoid: the actual window with the smalles total distance to the other windows in its 
cluster. 
- it starts with k-means ++ (10 restarts), then alternates assigning windows to the nearest centre and updating the 
medoids until the asignments stop changing.
- the crisis cluster is the one whose windows have the higher average S\&P variance
\subsubsection{Step 4}
- score each day
- on day d, the window ending d-1 is compared with both centres: $p_{\text{crisis}} = W_2(\text{calm})^2 / (W_2(\text{calm})^2 + W_2(\text{crisis})^2)$
- example: if the distance to calm is 2 and to crisis is 1, $p_{\text{crisis}} = 4 / (4 + 1) = 0.8$
- $p_{\text{crisis}} > 0.5$ exactly when the window is closer to the crisis centre, which is the paper's nearest centroid 
rule. So MVO and HRP switch exactly as the paper would classify
- it uses data through d-1 only, teh same timing as teh other two detectors.

\section{Regime Interpretation}
\subsection{MVO}
- it reads teh detector's daily csv and takes teh value for rebalance date d, which uses data through d-1
- before the detector's first value, $p_{\text{crisis}}$ is set to 0.5, meaning no information
- with \texttt{DETECTOR = "none"}, $p_{\text{crisis}}$ is 0, every model is its normal version
\subsection{MVO}
- on each rebalance date d, \texttt{in\_crisis(d)} checks $p_{\text{crisis}}(d) > 0.5$
- if yes, the weights are minimum-variance, if no, max-sharpe
- the choice is held until the next rebalance date. 
- it doesnt matter if it is 0.51 or 0.99, same effect
\subsection{HRP}
- on each rebalance date d, \texttt{in\_crisis(d)} checks $p_{\text{crisis}}(d) > 0.5$
- if yes, the weights are minimum-variance, if no, max-sharpe
- the choice is held until the next rebalance date. 
- it doesnt matter if it is 0.51 or 0.99, same effect
\subsection{XGB and RF}
- 




\section{Regimes}

\subsection{Changepoint Detection}
Methodology and calculation 1:1 to the paper


\subsection{Monthly Crisis Probability}

\begin{figure}[htbp]
    \centering
    \begin{tikzpicture}
        \renewcommand{\crisisshadecolor}{gray!15}%
        \renewcommand{\crisislinecolor}{black!45}%
        \pcrisispanel{changepoint}{Changepoint}{at={(0,0)}, xticklabel style={opacity=0}}
        \pcrisispanel{clustering}{Clustering}{at={(0,-\regimepanelheight)}, xticklabel style={opacity=0}}
        \pcrisispanel{hmm}{HMM}{at={(0,-2*\regimepanelheight)}}
    \end{tikzpicture}
    \caption{Monthly crisis probability $p_{\mathrm{crisis}}$ of the three regime detectors: changepoint detection (VIX, top), clustering (Wasserstein k-means, S\&P 500 and VIX, middle) and Hidden Markov Model (S\&P 500, bottom). Value on the first trading day of each month, held until the next month, as read by the monthly rebalanced portfolios; it uses data up to the previous day. Dashed: 0.5 threshold, above which the portfolios switch to their crisis allocation. Grey shading: crisis periods from the S\&P~500 peak until it is back at the peak; vertical lines: peak, trough and recovery.}
    \label{fig:pcrisis_monthly}
\end{figure}


\FloatBarrier
\section{Portfolio Performance}

\begin{table}[htbp]
    \centering
    \caption{Crisis overview: duration and maximum drawdown (\%) within each crisis; baselines, monthly rebalancing. Trading Days: S\&P~500 trading days from the peak to the trough. Maximum drawdown: from the S\&P~500 peak until the S\&P~500 is back at that peak. GSPC: S\&P~500 price index.}
    \label{tab:crisis_overview}
    \vspace{4pt}
    % row 0 of the file = duration (days), shown in the header; drawdowns stored as text
    % like -19.3%: read with % as a normal character, printed as they are
    \catcode`\%=12 \pgfplotstableread[col sep=comma]{\tableoverviewfile}\overview \catcode`\%=14
    \newcommand{\ovdays}[1]{\pgfplotstablegetelem{0}{#1}\of\overview\pgfplotsretval}
    % crisis name over two lines, bold
    \newcommand{\ovname}[2]{\textbf{\begin{tabular}[b]{@{}c@{}}#1\\#2\end{tabular}}}
    \footnotesize
    \setlength{\tabcolsep}{3pt}
    \pgfplotstabletypeset[
        skip rows between index={0}{1},
        every head row/.style={
            before row=\toprule,
            after row={\midrule
                       Trading Days & \ovdays{C1} & \ovdays{C2} & \ovdays{C3} & \ovdays{C4}
                       & \ovdays{C5} & \ovdays{C6} & \ovdays{C7}\\ \midrule
                       \textbf{Strategy} & \multicolumn{7}{c}{\textbf{Maximum Drawdown}}\\ \midrule}},
        every last row/.style={after row=\bottomrule},
        % rows counted after the skip: 0 GSPC, 1 MCap, 2 1/N, 3 HRP, 4 MVO, 5 LSTM, ...
        every row no 0/.style={after row=\midrule},             % GSPC | benchmarks
        every row no 2/.style={after row=\midrule},             % 1/N | HRP
        every row no 4/.style={after row=\midrule},             % MVO | LSTM
        string type, column type=c,
        columns/model/.style={column name={}, column type=l},
        columns/C1/.style={column name=\ovname{Asian}{Crisis}},
        columns/C2/.style={column name=\ovname{Dotcom}{Crisis}},
        columns/C3/.style={column name=\ovname{Financial}{Crisis}},
        columns/C4/.style={column name=\ovname{Monetary}{Policy Shock}},
        columns/C5/.style={column name=\ovname{COVID-19}{Crisis}},
        columns/C6/.style={column name=\ovname{Inflation \&}{Rate Hikes}},
        columns/C7/.style={column name=\ovname{Trade Policy}{Shock}},
    ]{\overview}
\end{table}

\begin{table}[htbp]
    \centering
    \caption{Full-period performance 1998--2025; baselines, monthly rebalancing. Sharpe ratio on excess returns over the 3-month T-bill, annualised. CAGR: compound annual growth rate. Cum.\ return: cumulative return over the period. Ulcer index: root mean square drawdown from the running high. Turnover p.a.: average annual turnover, without the initial purchase. GSPC: S\&P~500 price index (without dividends), held buy \& hold; the portfolios are total-return series.}
    \label{tab:full_period}
    \vspace{4pt}
    % cells stored as text (-56.8%, 0.25): read with % as a normal character, printed as they are
    \catcode`\%=12 \pgfplotstableread[col sep=comma]{\tablefullperiodfile}\fullperiod \catcode`\%=14
    \newcommand{\fpname}[2]{\textbf{\begin{tabular}[b]{@{}c@{}}#1\\#2\end{tabular}}}
    \footnotesize
    \setlength{\tabcolsep}{4pt}
    \pgfplotstabletypeset[
        every head row/.style={before row=\toprule, after row=\midrule},
        every last row/.style={after row=\bottomrule},
        % 0 GSPC, 1 MCap, 2 1/N, 3 HRP, 4 MVO, 5 LSTM, ...
        every row no 0/.style={after row=\midrule},             % GSPC | benchmarks
        every row no 2/.style={after row=\midrule},             % 1/N | HRP
        every row no 4/.style={after row=\midrule},             % MVO | LSTM
        string type, column type=c,
        columns/model/.style={column name=\textbf{Strategy}, column type=l},
        columns/sharpe_ratio/.style={column name=\fpname{Sharpe}{Ratio}},
        columns/cagr/.style={column name=\fpname{}{CAGR}},
        columns/cum_return/.style={column name=\fpname{Cum.}{Return}},
        columns/annual_volatility/.style={column name=\fpname{Annual}{Volatility}},
        columns/max_drawdown/.style={column name=\fpname{Max.}{Drawdown}},
        columns/ulcer_index/.style={column name=\fpname{Ulcer}{Index}},
        columns/turnover_pa/.style={column name=\fpname{Turnover}{p.a.}},
    ]{\fullperiod}
\end{table}

\begin{table}[htbp]
    \centering
    \caption{Average crisis performance: equal-weighted average over the seven crises; baselines, monthly rebalancing. Fixed S\&P~500 windows: full crisis from the S\&P~500 peak until it is back at that peak, peak to trough, and trough to recovery (back at the peak). Max.\ drawdown and Ulcer index over the full crisis; cumulative return over each of the three windows. Recovery: trading days from the portfolio's own trough until it is back at its value on the S\&P~500 peak date, searched up to the next crisis peak; ≥: not recovered in at least one crisis (counted up to that cut-off, a lower bound); MCap and MVO did not recover from the Dotcom Crisis before the peak of the Financial Crisis (2007-10-09), so their Dotcom Crisis recovery is counted up to that date. GSPC: S\&P~500 price index; its full-crisis return is about zero by construction of the windows.}
    \label{tab:crisis_average}
    \vspace{4pt}
    % cells stored as text (-31.9%, ≥377): read with % as a normal character, printed as they are
    \catcode`\%=12 \pgfplotstableread[col sep=comma]{\tablecrisisavgfile}\crisisavg \catcode`\%=14
    \newcommand{\caname}[2]{\textbf{\begin{tabular}[b]{@{}c@{}}#1\\#2\end{tabular}}}
    \footnotesize
    \setlength{\tabcolsep}{4pt}
    \pgfplotstabletypeset[
        every head row/.style={
            before row={\toprule & & & & \multicolumn{3}{c}{\textbf{Cumulative Return}}\\ \cmidrule(lr){5-7}},
            after row=\midrule},
        every last row/.style={after row=\bottomrule},
        % 0 GSPC, 1 MCap, 2 1/N, 3 HRP, 4 MVO, 5 LSTM, ...
        every row no 0/.style={after row=\midrule},             % GSPC | benchmarks
        every row no 2/.style={after row=\midrule},             % 1/N | HRP
        every row no 4/.style={after row=\midrule},             % MVO | LSTM
        string type, column type=c,
        columns/model/.style={column name=\textbf{Strategy}, column type=l},
        columns/avg_max_dd/.style={column name=\caname{Max.}{Drawdown}},
        columns/avg_ulcer/.style={column name=\caname{Ulcer}{Index}},
        columns/avg_recovery_days/.style={column name=\caname{Recovery}{(days)}},
        columns/avg_return_full/.style={column name=\caname{Full}{Crisis}},
        columns/avg_return_ptt/.style={column name=\caname{Peak to}{Trough}},
        columns/avg_return_ttp/.style={column name=\caname{Trough to}{Recovery}},
    ]{\crisisavg}
\end{table}

\begin{table}[htbp]
    \centering
    \caption{Regime detectors: changepoint (VIX), Hidden Markov Model (S\&P~500) and clustering (Wasserstein k-means, S\&P~500 and VIX), 1998--2025. Daily: the detector's own daily state. Monthly, quarterly, yearly: the state on the first trading day of the month, quarter or year, held until the next one, as the portfolios hold it. Episodes: number of switches into the crisis state; duration: their average length in trading days. Defensive: share of trading days in the crisis state. Entry lag: trading days from the S\&P~500 peak to the first crisis day (negative: already in the crisis state at the peak), over the crises entered before the trough. Exit lag: trading days from the S\&P~500 trough to the first calm day, over the crises in which the detector switches back. In brackets: number of the seven crises behind the average lag. Recall and precision per day against the crisis label S\&P~500 peak to trough.}
    \label{tab:detectors}
    \vspace{4pt}
    % cells stored as text (48.7%, -74.3 (6)): read with % as a normal character, printed as they are
    \catcode`\%=12 \pgfplotstableread[col sep=comma]{\tabledetectorsfile}\detectors \catcode`\%=14
    \newcommand{\dtname}[2]{\textbf{\begin{tabular}[b]{@{}c@{}}#1\\#2\end{tabular}}}
    \footnotesize
    \setlength{\tabcolsep}{4pt}
    \pgfplotstabletypeset[
        columns={detector,n_episodes,avg_episode_days,defensive_share,avg_entry_lag,avg_exit_lag,recall,precision},
        every head row/.style={before row=\toprule, after row=\midrule},
        every last row/.style={after row=\bottomrule},
        % the frequency column is not shown: a header row above each block of three detectors
        every row no 0/.style={before row={\multicolumn{8}{l}{\textit{Daily}}\\}},
        every row no 3/.style={before row={\midrule\multicolumn{8}{l}{\textit{Monthly}}\\}},
        every row no 6/.style={before row={\midrule\multicolumn{8}{l}{\textit{Quarterly}}\\}},
        every row no 9/.style={before row={\midrule\multicolumn{8}{l}{\textit{Yearly}}\\}},
        string type, column type=c,
        columns/detector/.style={column name=\textbf{Detector}, column type=l},
        columns/n_episodes/.style={column name=\dtname{}{Episodes}},
        columns/avg_episode_days/.style={column name=\dtname{Duration}{(days)}},
        columns/defensive_share/.style={column name=\dtname{}{Defensive}},
        columns/avg_entry_lag/.style={column name=\dtname{Entry Lag}{(days)}},
        columns/avg_exit_lag/.style={column name=\dtname{Exit Lag}{(days)}},
        columns/recall/.style={column name=\dtname{}{Recall}},
        columns/precision/.style={column name=\dtname{}{Precision}},
    ]{\detectors}
\end{table}

\begin{table}[htbp]
    \centering
    \caption{Detector influence: average over HRP, MVO, LSTM, RF and XGB, monthly rebalancing. Baseline: without regime detection. Cash: T-bills while the detector signals a crisis. Alternative: minimum variance (HRP, MVO) or HRP (LSTM, RF, XGB) while the detector signals a crisis. Each value is first averaged over the seven crises per model, then over the five models, all with equal weight. Cumulative return over the fixed S\&P~500 windows peak to trough, trough to recovery and peak to recovery (full crisis); maximum drawdown and Ulcer index over the full crisis. Sharpe ratios pooled per model, then averaged over the five models: in crisis over the daily returns of all seven full crises together, outside crisis over all other trading days 1998--2025; annualised, in excess of the T-bill rate. Crisis Sharpe ratios pool the daily returns of all seven crises, so they are weighted by crisis length (the Dotcom and Financial crises account for about 77\% of crisis days); all other columns weight each crisis equally.}
    \label{tab:detector_influence}
    \vspace{4pt}
    % cells stored as text (-29.8%, 0.02): read with % as a normal character, printed as they are
    \catcode`\%=12 \pgfplotstableread[col sep=comma]{\tabledetectorinfluencefile}\influence \catcode`\%=14
    \newcommand{\diname}[2]{\textbf{\begin{tabular}[b]{@{}c@{}}#1\\#2\end{tabular}}}
    \footnotesize
    \setlength{\tabcolsep}{4pt}
    \pgfplotstabletypeset[
        every head row/.style={
            before row={\toprule & \multicolumn{3}{c}{\textbf{Cumulative Return}} & & & \multicolumn{2}{c}{\textbf{Sharpe Ratio}}\\
                        \cmidrule(lr){2-4} \cmidrule(lr){7-8}},
            after row=\midrule},
        every last row/.style={after row=\bottomrule},
        % 0 Baseline, 1-2 Changepoint, 3-4 Clustering, 5-6 HMM
        every row no 0/.style={after row=\midrule},
        every row no 2/.style={after row=\midrule},
        every row no 4/.style={after row=\midrule},
        string type, column type=c,
        columns/strategy/.style={column name=\textbf{Strategy}, column type=l},
        columns/ret_ptt/.style={column name=\diname{Peak to}{Trough}},
        columns/ret_ttp/.style={column name=\diname{Trough to}{Recovery}},
        columns/ret_full/.style={column name=\diname{Full}{Crisis}},
        columns/max_dd/.style={column name=\diname{Max.}{Drawdown}},
        columns/ulcer/.style={column name=\diname{Ulcer}{Index}},
        columns/sharpe_crisis/.style={column name=\diname{In}{Crisis}},
        columns/sharpe_outside/.style={column name=\diname{Outside}{Crisis}},
    ]{\influence}
\end{table}

\begin{figure}[htbp]
    \centering
    \begin{tikzpicture}
        \renewcommand{\crisisshadecolor}{gray!15}%
        \renewcommand{\crisislinecolor}{black!45}%
        \valuepanel{hrp}{HRP}{at={(0,0)}, xticklabel style={opacity=0},
            legend columns=4, legend cell align=left,
            legend style={at={(0.5,1.12)}, anchor=south, font=\footnotesize, draw=none,
                          /tikz/every even column/.append style={column sep=6pt}},
            legend entries={Baseline, Changepoint Cash, Clustering Cash, HMM Cash,
                            \mbox{}, Changepoint Alternative, Clustering Alternative, HMM Alternative}}
            {\crisislabels}
        \valuepanel{lstm}{LSTM}{at={(0,-\valuepanelheight)}}{}
    \end{tikzpicture}
    \caption{Value of 1 invested in January 1998 (log scale) of the seven monthly rebalanced portfolios of HRP (top) and LSTM (bottom): baseline without regime detection, and per detector the cash variant (dashed, T-bills while the detector signals a crisis) and the alternative variant (solid; HRP: minimum variance, LSTM: HRP while the detector signals a crisis). Same colours in both panels. Grey shading: crises C1--C7 from the S\&P~500 peak until it is back at the peak; vertical line: S\&P~500 trough.}
    \label{fig:value_hrp_lstm}
\end{figure}

\begin{figure}[p]
    \centering
    \begin{tikzpicture}
        \begin{groupplot}[
            group style={group size=2 by 4, horizontal sep=1.4cm, vertical sep=1.9cm,
                         ylabels at=edge left},
            width=0.45\textwidth, height=3.7cm, scale only axis=false,
            ylabel={Cum.\ log return},
            title style={font=\small\bfseries, yshift=-2pt},
            tick label style={font=\scriptsize},
            xticklabel style={rotate=40, anchor=north east, inner sep=1pt},
            yticklabel style={/pgf/number format/fixed},
            grid=major, grid style={gray!20},
        ]
            % {k}{title}{1st of peak month}{peak}{trough}{even}{end of even month}
            \crisispanel{1}{C1: Asian Crisis (1998)}{1998-07-01}{1998-07-17}{1998-08-31}{1998-11-23}{1998-11-30}
            \crisispanel{2}{C2: Dotcom Crisis (2000)}{2000-03-01}{2000-03-24}{2002-10-09}{2007-05-30}{2007-05-31}
            \crisispanel{3}{C3: Financial Crisis (2007)}{2007-10-01}{2007-10-09}{2009-03-09}{2013-03-28}{2013-03-31}
            \crisispanel{4}{C4: Monetary Policy Shock (2018)}{2018-09-01}{2018-09-20}{2018-12-24}{2019-04-23}{2019-04-30}
            \crisispanel{5}{C5: COVID-19 Crisis (2020)}{2020-02-01}{2020-02-19}{2020-03-23}{2020-08-18}{2020-08-31}
            \crisispanel{6}{C6: Inflation \& Rate Hikes (2022)}{2022-01-01}{2022-01-03}{2022-10-12}{2024-01-19}{2024-01-31}
            \crisispanel{7}{C7: Trade Policy Shock (2025)}{2025-02-01}{2025-02-19}{2025-04-08}{2025-06-27}{2025-06-30}
            % eighth slot: the legend
            \nextgroupplot[hide axis, xmin=0, xmax=1, ymin=0, ymax=1, ylabel={},
                           legend columns=2, legend cell align=left,
                           legend style={at={(0.5,0.5)}, anchor=center, draw=none, font=\footnotesize,
                                         row sep=2pt, /tikz/every even column/.append style={column sep=8pt}}]
                \addlegendimage{black, line width=1.1pt}        \addlegendentry{LSTM Baseline}
                \addlegendimage{gray!60, line width=0.5pt}      \addlegendentry{MCap}
                \addlegendimage{regcash=regCPcash}              \addlegendentry{Changepoint Cash}
                \addlegendimage{regalt=regCP}                   \addlegendentry{Changepoint Alternative}
                \addlegendimage{regcash=regCLcash}              \addlegendentry{Clustering Cash}
                \addlegendimage{regalt=regCL}                   \addlegendentry{Clustering Alternative}
                \addlegendimage{regcash=regHMMcash}             \addlegendentry{HMM Cash}
                \addlegendimage{regalt=regHMM}                  \addlegendentry{HMM Alternative}
                \addlegendimage{black!60, dashed, line width=0.5pt} \addlegendentry{S\&P~500 trough}
        \end{groupplot}
    \end{tikzpicture}
    \caption{Cumulative log return of the seven monthly rebalanced LSTM portfolios around each crisis, relative to the S\&P~500 peak (set to 0), weekly values; each panel spans the full calendar months from the month of the peak to the month in which the S\&P~500 is back at the peak, the crisis itself (peak to recovery) shaded light grey: baseline without regime detection (thick black), and per detector the cash variant (dashed, T-bills while the detector signals a crisis) and the alternative variant (solid, HRP while the detector signals a crisis); market-cap benchmark (MCap, thin grey) for orientation. Dashed vertical line: S\&P~500 trough, separating peak to trough and trough to recovery. Date ticks: peak, trough, recovery. Each panel has its own y scale, as the crises differ greatly in depth and length.}
    \label{fig:crises_lstm}
\end{figure}

\FloatBarrier

\subsection{Crisis Excess Returns: Peak to Trough vs.\ Trough to Recovery}

\begin{figure}[htbp]
    \centering
    \begin{tikzpicture}
        \begin{axis}[
            width=\textwidth,
            axis equal image,      % same scale on both axes: the net line is at 45 degrees
            xmin=-4, xmax=13, ymin=-5.5, ymax=6.5,
            xtick={-4,-2,...,12}, ytick={-4,-2,...,6},
            xlabel={Excess return vs.\ MCap monthly, peak to trough (\%)},
            ylabel={Excess return vs.\ MCap monthly, trough to recovery (\%)},
            grid=major, grid style={gray!20},
            tick label style={font=\small},
            table/col sep=comma,
            % dashed zero lines = benchmark level
            extra x ticks={0}, extra y ticks={0},
            extra tick style={grid=major, grid style={gray, dashed}, tick label style={opacity=0}},
            % legend below the plot: models (colour) / rebalancing (marker)
            legend style={at={(0.5,-0.14)}, anchor=north, legend columns=7,
                          font=\footnotesize, draw=gray!50,
                          /tikz/every even column/.append style={column sep=0.8em}},
            legend cell align=left,
            % one colour per model (key column), marker per frequency (plot)
            scatter/classes={
                mcap={modMCap}, equal={modEqual}, hrp={modHRP}, mvo={modMVO},
                lstm={modLSTM}, rf={modRF}, xgb={modXGB}},
        ]
            % net ~ 0 over the whole crisis: ptt + ttp = 0 (over the visible range)
            \addplot[forget plot, no marks, black!60, densely dotted, thick, domain=-4:5.5, samples=2]
                {-x} node[pos=0.12, sloped, above, font=\footnotesize] {net $\approx 0$};
            \node[font=\footnotesize, black!60, anchor=south west] at (axis cs:-3.9,-5.4) {net loss};
            \node[font=\footnotesize, black!60, anchor=north east] at (axis cs:12.9,6.4) {net gain};
            % arrows monthly -> quarterly -> yearly of the same model; row 0 (MCap,
            % monthly only, quarterly / yearly nan) is skipped
            \addplot[forget plot, gray!70, -stealth, line width=0.6pt, skip coords between index={0}{1},
                     quiver={u=\thisrow{ptt_quarterly}-\thisrow{ptt_monthly},
                             v=\thisrow{ttp_quarterly}-\thisrow{ttp_monthly},
                             every arrow/.append style={shorten <=2.5pt, shorten >=2.5pt}}]
                table[x=ptt_monthly, y=ttp_monthly] {\scatterfile};
            \addplot[forget plot, gray!70, -stealth, line width=0.6pt, skip coords between index={0}{1},
                     quiver={u=\thisrow{ptt_yearly}-\thisrow{ptt_quarterly},
                             v=\thisrow{ttp_yearly}-\thisrow{ttp_quarterly},
                             every arrow/.append style={shorten <=2.5pt, shorten >=2.5pt}}]
                table[x=ptt_quarterly, y=ttp_quarterly] {\scatterfile};
            % points: monthly circle, quarterly square, yearly triangle
            \addplot[forget plot, scatter, only marks, mark=*, mark size=2.3pt, point meta=explicit symbolic]
                table[x=ptt_monthly, y=ttp_monthly, meta=key] {\scatterfile};
            \addplot[forget plot, scatter, only marks, mark=square*, mark size=2pt, point meta=explicit symbolic]
                table[x=ptt_quarterly, y=ttp_quarterly, meta=key] {\scatterfile};
            \addplot[forget plot, scatter, only marks, mark=triangle*, mark size=2.6pt, point meta=explicit symbolic]
                table[x=ptt_yearly, y=ttp_yearly, meta=key] {\scatterfile};
            % model names at the monthly point; side from column label_pos
            \addplot[
                forget plot, only marks, mark=none,
                visualization depends on={value \thisrow{model} \as \modelname},
                visualization depends on={value \thisrow{label_pos} \as \labelpos},
                nodes near coords={\expandafter\detokenize\expandafter{\modelname}},
                nodes near coords style={font=\footnotesize, anchor=\labelanchor{\labelpos}, inner sep=3pt},
            ] table[x=ptt_monthly, y=ttp_monthly] {\scatterfile};
            % legend
            \addlegendimage{only marks, mark=*, modMCap}  \addlegendentry{MCap}
            \addlegendimage{only marks, mark=*, modEqual} \addlegendentry{1/N}
            \addlegendimage{only marks, mark=*, modHRP}   \addlegendentry{HRP}
            \addlegendimage{only marks, mark=*, modMVO}   \addlegendentry{MVO}
            \addlegendimage{only marks, mark=*, modLSTM}  \addlegendentry{LSTM}
            \addlegendimage{only marks, mark=*, modRF}    \addlegendentry{RF}
            \addlegendimage{only marks, mark=*, modXGB}   \addlegendentry{XGB}
            \addlegendimage{only marks, mark=*, black}            \addlegendentry{Monthly}
            \addlegendimage{only marks, mark=square*, black}      \addlegendentry{Quarterly}
            \addlegendimage{only marks, mark=triangle*, black}    \addlegendentry{Yearly}
        \end{axis}
    \end{tikzpicture}
    \caption{Average crisis excess return over the market-cap benchmark (MCap, monthly rebalancing) per model and rebalancing frequency: peak to trough ($x$) vs.\ trough to recovery ($y$), equal-weighted average over the seven crises, fixed S\&P~500 windows. 19 portfolios: MCap monthly (at the origin) and the baselines of 1/N, HRP, MVO, LSTM, RF and XGB at monthly (circle), quarterly (square) and yearly (triangle) rebalancing; arrows connect the same model from monthly to quarterly to yearly. Excess return $(1+R)/(1+R_{\mathrm{MCap}})-1$ per crisis and phase. The dotted line marks $x + y = 0$: portfolios above it gain over the whole crisis (peak to recovery), portfolios below lose. The line is approximate, since the two phases compound rather than add.}
    \label{fig:scatter_crisis_excess}
\end{figure}
\FloatBarrier


\subsection{Detection Quality}

\FloatBarrier




\newpage
\section{References}

Nystup Peter, Madsen Hendrik, Lindström Erik, 2017. Dynamic portfolio optimization across hidden market regimes.\\ 
Used for:\\
    - Transaction costs 10bps\\
    - Hidden Markov Model\\




    



\end{document}